% answers_short.tex —— 2012 年·广东卷·文科 · 简短答案速查
% 由 tools/make_short_answers.py 自 answers.tex 的 \gksolhead/\gkanswer 提取，
% 并与 2012/2012_数学真题_广东文_答案要点.md 逐题比对后生成，勿手改（改源 MD 后重跑）。

\begin{gknotice}
\textbf{使用说明}：本节只给最终答案便于快速核对，完整推导见解析版。答案为编写者独立推导，\textbf{非官方标准答案}。
\end{gknotice}

\gkpart{选择题}
\begin{center}\begin{tabular}{ccccccccccc}
\toprule
题号 & 1 & 2 & 3 & 4 & 5 & 6 & 7 & 8 & 9 & 10 \\
\midrule
答案 & D & A & A & D & C & B & C & B & C & D \\
\bottomrule
\end{tabular}\end{center}

\gkpart{填空题}
\begin{itemize}[leftmargin=2.6em]
  \item[\textbf{11.}] $[-1,0)\cup(0,+\infty)$。需 $x+1\geqslant0$ 且 $x\neq0$。
  \item[\textbf{12.}] $\dfrac14$。$a_3^2=a_2a_4=\dfrac12$，$a_1a_5=a_3^2$，故原式 $=\left(\dfrac12\right)^2=\dfrac14$。
  \item[\textbf{13.}] $1,1,3,3$。四数和为 $8$、中间两数和为 $4$，逐一检验方差为 $1$ 的组合，仅此一组。
  \item[\textbf{14.}] $(2,1)$。$C_1$ 是圆 $x^2+y^2=5$ 在第一象限的四分之一圆弧，$C_2$ 是直线 $y=x-1$，联立得 $x=2$ 或 $x=-1$，舍去后者。
  \item[\textbf{15.}] $\sqrt{mn}$。由弦切角定理及 $\angle PBA=\angle DBA$ 得 $\angle DBA=\angle BCA$，故 $\triangle ABD\backsim\triangle ACB$，$AB^2=AD\cdot AC=mn$。
\end{itemize}

\gkpart{解答题}
\begin{itemize}[leftmargin=2.6em]
  \item[\textbf{16.}] (1) $A=2$。$f\left(\dfrac\pi3\right)=A\cos\dfrac\pi4=\dfrac{\sqrt2}{2}A=\sqrt2$；(2) $-\dfrac{13}{85}$。由条件得 $\sin\alpha=\dfrac{15}{17}$，$\cos\alpha=\dfrac{8}{17}$，$\cos\beta=\dfrac45$，$\sin\beta=\dfrac35$，再用两角和的余弦公式
  \item[\textbf{17.}] (1) $a=0.005$。各组矩形高之和乘组距 $10$ 等于 $1$，即 $10(2a+0.09)=1$；(2) $73$。各组取中点乘以对应频率再求和；(3) $10$。语文各段人数为 $5,40,30,20$，按比例得数学各段人数 $5,20,40,25$，合计 $90$，故 $[50,90)$ 之外有 $100-90=10$ 人
  \item[\textbf{18.}] (1) 见解析版；(2) $\dfrac{\sqrt2}{12}$。$E$ 为 $PB$ 中点，故 $V_{E-BCF}=\dfrac12V_{P-BCF}$，而 $S_{\triangle BCF}=\dfrac12\cdot FC\cdot AD=\dfrac{\sqrt2}{2}$，$PH=1$；(3) 见解析版
  \item[\textbf{19.}] (1) $a_1=1$。$n=1$ 时 $T_1=S_1=a_1$，代入得 $a_1=2a_1-1$；(2) $a_n=3\cdot2^{n-1}-2$。由 $T_n-T_{n-1}=S_n$ 得 $S_n=2S_{n-1}+2n-1$，解得 $S_n=3\cdot2^n-2n-3$，再作差
  \item[\textbf{20.}] (1) $\dfrac{x^2}{2}+y^2=1$。$P(0,1)$ 给 $b=1$，$c=1$ 给 $a^2=b^2+c^2=2$；(2) $x-\sqrt2y+2=0$ 或 $x+\sqrt2y+2=0$。设 $l:y=kx+m$，与抛物线相切给 $km=1$，与椭圆相切给 $m^2=1+2k^2$，联立得 $k^2=\dfrac12$
  \item[\textbf{21.}] (1) $0<a<\dfrac13$ 时 $D=(0,x_1)\cup(x_2,+\infty)$，其中 $x_{1,2}=\dfrac{3(1+a)\mp\sqrt{9a^2-30a+9}}4$；$a=\dfrac13$ 时 $D=(0,1)\cup(1,+\infty)$；$\dfrac13<a<1$ 时 $D=(0,+\infty)$。按判别式 $\Delta=3(3a-1)(a-3)$ 的符号分三种情形；(2) $0<a<\dfrac13$ 时只有极大值点 $x=a$；$a=\dfrac13$ 时只有极大值点 $x=\dfrac13$；$\dfrac13<a<1$ 时极大值点 $x=a$、极小值点 $x=1$。由 $f'(x)=6(x-1)(x-a)$ 及 $a$、$1$ 是否落在 $D$ 内判定
\end{itemize}
